%TCIDATA{Version=5.50.0.2890}
%TCIDATA{LaTeXparent=0,0,NoticeDKE.tex}

The determination of the electron distribution function has a crucial
importance in the tokamak plasma physics, since the toroidal current density
profile that is mainly driven by electrons is intimately linked to the
magnetic equilibrium and confinement performances \cite{pey03}. Therefore,
accurate and realistic calculations must be carried out, with the additional
requirement of an optimized numerical approach, in order to reduce as much
as possible both memory storage and computer time consumptions. The latter
point is especially important, since kinetic calculations must be incorporated in a
chain of codes for self-consistent determination of all plasma properties
\cite{bas03}.

In this document, an extensive presentation of the fast solver for the
linearized electron drift kinetic is presented. This is a completely new
tool based on previous numerical developments \cite{sch99}, that is designed
for realistic calculations of the electron distribution function in the
plasma region where the weak collision banana regime holds. It incorporates
the major physical ingredients that must be taken into account for
describing the corresponding physics in a fusion reactor, namely
relativistic corrections, trapped particle effects, arbitrary magnetic
equilibrium for high $\beta _{p}$ regimes. For this purpose, both zero and
first order kinetic equations with respect to the small drift approximation
are solved, which allows to determine self-consistently boostrap current
with any type of external current source (RF, Ohmic,...) at any point of the
momentum space, and not only at the trapped-passing boundary as done
in a previous attempt \cite{wes96}. Basically, the
code gives access to the neoclassical physics dominated by collisions
between charged particles, for non-Maxwellian electron distribution
functions. Therefore, it is particularly well suited for accurate current
drive estimates in advanced tokamak regimes, including ITER, where locally,
bootstrap current may strongly interplay with external current sources
(ITB, edge pedestal in H-mode\ldots)

Besides these physical properties, the code offer also the possibility to
incorporate any type of fast electron radial transport (collision,
turbulence or wave induced), which may be a key ingredient for the local
control of plasma properties \cite{pey93}. Written in a fully conservative
form, the code naturally conserve the electron density, but also momentum
for the current drive problem, keeping first order term of the Legendre
polynomial expansion of the Beliaev-Budker collision operator \cite{bra89}.
As usual, several useful moments of the electron distribution function are
calculated, namely the current density, the absorbed power, the fraction of
trapped electrons, the magnetic ripple losses \cite{ju02} and the
bremsstrahlung emission \cite{pey99}.

Advanced numerical techniques have been used, so that memory storage
requirement can be strongly reduced, while keeping fast convergence rate.
For this purpose the electron Drift Kinetic equation is solved by the
standard finite difference technique, which has proven so far to be the
fastest numerical approach among all possible alternative techniques.
Furthermore, this method is particularly well suited when large
discontinuities of the diffusion or convection rates have to be considered,
a case that occurs frequently when kinetic and ray-tracing calculations are
coupled.

Since in most cases, the steady-state solution is seeked with respect to the
largest time scale (collision or fast electron radial
transport)\footnote{The energy transport time scale is usually on
order on magnitude larger than the largest characteristic time for
current drive calculations, except in tokamaks of small size, where
time ordering here considered in the model must be likely revisited}, the appropriate
technique is the well known upwind time differencing, corresponding to the
fully implicit time scheme, whose characteristic is to be almost
unconditionally stable with respect to the time step value $\Delta t$. Nevertheless, the
code offers also the possibility to investigate time dependent problems,
with the usual Crank-Nicholson time differencing, which enables accurate
time evolution.

The bounce-averaged Drift Kinetic equation is basically a $3-D$ problem, $%
2-D $ in momentum space (slowing down, pitch-angle) et $1-D$ in
configuration space (radial dimension). Up to now, in order to reduce memory
storage requirements, the numerical time scheme was based on the operator
splitting technique, where both momentum and spatial dynamics evolved
separately. If this approach turns out to be very fruitful, it has the
drawback to slow down considerably the convergence towards the steady state
solution, since only small time steps may be used for numericaly stable
convergence. Therefore, the advantage to use fully implicit time scheme for
each sub-space in hindered by this strong limitation, especially when radial
transport of fast electrons must be taken into account. In order to avoid
this problem, a fully implicit time scheme is considered, where both
momentum and spatial dynamics are simultaneaously considered, so that no
limitation occurs on the time step, which may be several order of magnitude
larger than the collision reference time. However, this method requires a
new technique for matrix inversion, in order to keep memory storage at an
acceptable level. Indeed, with usual mesh sizes, the standard $LU$ matrix
factorization techniques does not hold anymore since matrices requirement
may reach several giga-Bytes. An alternative approach is therefore
absolutely necessary.

This critical point has been addressed by using advanced inversion
techniques, based on incomplete $LU$ factorization with drop tolerance.
Since most of the off-diagonal coefficients of the matrices $L$ and $U$ are
very small, one may take advantage to remove them so that the sparsity of
the matrices can be greatly enhanced. Memory storage requirements can be
therefore reduced drastically by one or two orders of magnitude with this
pruning method, depending upon the initial matrix preconditioning, while
only coefficients that are relevant of the physics problem here addressed
are kept. Furthermore, computer time consumption can be also reduced, since
the number of non-zero coefficients is considerably reduced. This method is
similar to the strongly implicit method used for factorizing nine diagonal
matrices \cite{pey98}, except that in this case, no restriction takes place
regarding the number of diagonals. However, to take full benefit of this
approach, the non-zero elements of the matrix which is inverted must lie
predominantly along diagonals. Therefore, it may be applied for solving the
zero and first order bounce-averaged Fokker-Planck equations, whose
structures are well suited for this purpose, though coefficients arrangement
can be complex, owing to the radial dependence of the internal
trapped-passing boundary in momentum space, especialy when transport in
configuration space has to be considered.

This approach has been very successfuly implemented for the electron Drift
Kinetic problem in tokamaks, using the MatLab\ language, which provides a
built-in function for incomplete LU factorization with drop tolerance, and
several very efficient iterative inversion tools, like the Conjugate
Gradient Squared method for solving the system of linear equations. It
is important to recall that this method is also available in FORTRAN
programming language\footnote{Useful informations are available on the website of Pr. Yousef Saad
at the following internet address http://www-users.cs.umn.edu/\symbol{126}saad/}, under the package name
SPARSEKIT that has also parallel processing capabilities \cite{saa03}.
Moreover, the very compact MatLab\ programming syntax allows to design the
code structure in an original way, using multidimensional objects that
describe simultaneously momentum and configuration space dynamics, but also
wave-particle interaction. This makes the code particularly robust and easy
to maintain.

In the document, the physics and numerics issues of the code are detailed,
and an extensive discussion of the underlying assumptions is presented. A
specific attention is paid to derive matrix coefficients in a fully
consistent manner, a crucial issue especially for an accurate and robust
estimate of the current drive efficiency for the various methods used in tokamaks. Some examples
are shown to illustrate code performances, though still numerous
possibilities remain to be investigated but are beyond the purpose of
this document.

Aside from present day code capabilities, it is important to notice that the
new numerical approach, here used, gives access to new physics domains that
have never been studied accurately like wave-induce radial transport
\cite{kup91}. Furthermore, since the algorithm used is fast and stable, possible
extensions to $4-D$ problems may be foreseen like in the plateau collision
regime (current drive at the very plasma edge), as well as studies of the
difficult problem of electrons that are locally trapped at different spatial
positions, like in stellarator. In addition the code may be extended quite
in a straightforward manner to the multi-species problem, taking into
account for example of the non-linear damping of the $\alpha$-particles
produced by fusion reactions on the electron population. However, the ion
physics requires to perform orbit-averaging instead of
bounce-averaging, because of the large banana width of some particles, a
challenging issue for kinetic solver based on a finite difference technique.
Such a requirement is crucial for describing torque induced by
waves. Nevertheless, beside this difficulty, the code is already fully
designed to take benefit from parallel processing, if the dynamics of
various species must be studied. In particular, non-uniform momentum
and pitch-angle grids are already implemented, so that refined calculations can be
performed for the ions at low velocity, while accurate ones up to
relativistic energies may be considered for the electrons.
